\chapter{Introduction}
\label{chap:intro}

As modern software systems become increasingly complex, their attack surface expands correspondingly, elevating software vulnerabilities to a critical security concern. The consequences of unpatched software flaws can range from isolated service disruption to systemic data breaches with global impact. Automated vulnerability discovery has therefore transitioned from an academic pursuit to a fundamental necessity within the software development life cycle. While techniques such as static analysis \cite{nachtigall2022large} and symbolic execution \cite{baldoni2018survey, king1976symbolic} provide valuable insights and high precision, practical constraints regarding scalability and path explosion can make them less effective in large systems. In contrast, dynamic analysis --- and more specifically, fuzz testing or "fuzzing" --- has emerged as the remarkably pragmatic and dominant approach for identifying flaws in compiled binaries.

Coverage-Guided Fuzzing (CGF) has established itself as the \textit{de facto} standard for automated software testing. Evolving from early blind random-input generators \cite{miller1990empirical} to highly sophisticated feedback-driven engines \cite{fioraldi2020afl++, honggfuzz, libfuzzer}, fuzzing is now deployed on an industrial scale. It is responsible for discovering thousands of critical vulnerabilities across diverse domains, including general open-source applications \cite{ossfuzz}, Internet of Things (IoT) devices \cite{chen2018iotfuzzer, zheng2019firm}, operating system kernels \cite{syzkaller, xu2020krace}, and enterprise-level database systems \cite{jiang2023dynsql, liang2024wingfuzz, wang2021reinforcement}. Modern fuzzers rely on the continuous aggregation of execution telemetry to guide the generation of new test cases. By monitoring the control-flow graph (CFG) transitions triggered by mutated inputs, the fuzzer can actively prioritise inputs that discover previously unseen code paths, effectively exploring the target application's state space.

Despite its undeniable success, the efficacy of any CGF engine is fundamentally bound by the quality and granularity of the feedback guiding it \cite{wang2019sensitive}. The fuzzer relies on coverage metrics as a proxy for program state exploration. For this reason, several new techniques have been developed recently to overcome the limitations of CGF, such as hard-to-bypass constraints \cite{aschermann2019redqueen, poeplau2020symbolic, stephens2016driller, yun2018qsym} and the high number of invalid test cases generated by mutation \cite{aschermann2019nautilus, blazytko2019grimoire, fioraldi2020weizz, padhye2019semantic}. These modifications are all designed to maximise code coverage, as exploring a greater proportion of the codebase is empirically proven to yield a higher rate of flaw discovery.

Modern fuzzers predominantly rely on edge coverage \cite{manes2019art} rather than simple basic block coverage. By tracking the specific control-flow transitions between blocks, edge coverage provides the fuzzer with more granular feedback on the program's internal state. This is traditionally implemented by assigning a unique identifier to each branch transition \cite{zalewski2014afltech}, allowing the engine to quickly record and recognise newly discovered execution paths during the fuzzing loop.

Indirect branches pose a uniquely complex challenge within this tracking mechanism. Unlike direct branches, where the target address is statically known, an indirect branch --- such as a function pointer invocation, a virtual method dispatch in C++, or a complex switch-case statement --- dynamically resolves its destination at runtime based on register states or memory contents. Consequently, a single indirect call site can map to numerous diverse targets, significantly increasing the number of edge transitions that must be tracked.

To overcome this limitation, this thesis proposes and evaluates a new architectural enhancement for AFL++: the introduction of a secondary coverage map dedicated exclusively to tracking the targets of indirect control-flow transfers. We hypothesise that indirect branch tracking might provide the fuzzer with a higher-fidelity representation of the program's dynamic execution state, thereby improving overall path discovery.

The core tenets of the proposed architecture are:
\begin{itemize}
    \item \textit{Isolation:} Indirect branch data is stored in an independent shared memory region, preventing interference with the primary edge coverage bitmap and preserving standard edge discovery.

    \item \textit{Collision-Free Identification:} The architecture employed guarantees that every indirect branch source in the binary is mapped to a unique slot.

    \item \textit{Probabilistic Tracking:} To maintain strict memory and performance bounds, each indirect source slot acts as a Bloom filter containing already-seen transitions. 

    \item \textit{Dynamic Sizing:} The required capacity of the secondary map is computed iteratively at startup by the instrumented binary and communicated to the fuzzer, ensuring optimal memory utilisation.
\end{itemize}

We implemented this architecture by extending the LLVM compiler infrastructure to inject tracking logic at the site of every indirect jump and call, modifying the compiler-rt libraries to handle runtime initialisation, and altering the fuzzer's core orchestration to accommodate dynamic memory reallocation and dual-map seed scoring.

Following the implementation, an empirical evaluation was conducted using a diverse set of real-world benchmarks to assess the impact of the indirect coverage map on path discovery and fuzzing throughput. The results of this evaluation indicate that, while the proposed architecture functions mechanically as intended, it did not yield the anticipated breakthrough in overall coverage growth. Across all tested benchmarks, the baseline AFL++ consistently performed either on par with or slightly better than the modified architecture.

This thesis explicitly details these negative results and provides insight into the possible structural and theoretical limitations that caused them. We establish that the lack of improvement is primarily attributable to redundancy with the primary edge map. Furthermore, we identify limitations within the standard mutation engine, which struggles to generate the precise data structures (e.g., valid pointers) necessary to satisfy complex indirect dispatch constraints without the aid of specialised techniques \cite{aschermann2019redqueen}. Additional advanced taint-tracking approaches \cite{wang2010taintscope} might be required to fully leverage indirect paths.

\section{Contributions}

The primary contributions of this thesis are:
\begin{itemize}
    \item A comprehensive analysis of the limitations of traditional, fixed-size shared memory bitmaps in tracking indirect control-flow within Coverage-Guided Fuzzing.

    \item The design and implementation of an isolated, dynamically sized secondary coverage map for AFL++, providing collision-free tracking of indirect branch sources.

    \item An empirical evaluation of the proposed architecture against standard baselines, accompanied by a critical analysis of the structural limitations.
\end{itemize}

\section{Thesis Structure}

In Chapter \ref{ch:bg}, we provide the necessary background on automated software testing, detailing the terminology and a general description of fuzz testing. In Chapter \ref{ch:afl}, we dissect the architecture of AFL++. In Chapter \ref{ch:arch}, we present the theoretical design of the proposed architecture. In Chapter \ref{ch:impl}, we detail the concrete engineering modifications made to the AFL++ codebase. In Chapter \ref{ch:eval}, we present the results and analyse the observed outcomes. The thesis ends with Chapter \ref{ch:concl}, containing conclusions and the discussion for possible future directions.